\scanpage{8}
% archive sheet number: 7
% lines 1–2
\unsure{On|[On|Du} trouve en tous cas, sous la seule hyp. ($\alpha$) de la remarque, les conclusions :

% lines 3–6
$B\unsure{\leadsto|\sim} B\otimes_{A_X}\underline O_X$ est une équivalence de la catégorie des $A$-Alg. de présentation finie avec la catégorie des $\underline O_X$-Alg. de \uline{présentation} \uline{globale} finie.

% lines 7–11
Si $B$ \unsure{et|est} \deleted{\unsure{une}} $\underline B$ se correspondent, Alors $M\leadsto M\otimes_{B_X}\underline B$ est une équiv. de la catégorie des $B$-Modules de présentation finie, et la catégorie des $\underline B$-Modules des présentation \uline{globale} finie.

% lines 12–13
\uline{Question} Soit $\underline B$ de \inserted{Alg.} prés. finie sur $X$, le noyau $J$ de $\underline O_X\unsure{-|\to}B$ est-il cohérent ??

% lines 14–18
C'est une question de pure Algèbre comm., peut s'énoncer : si $A$ est un anneau cohérent, \unsure{et|\deleted{$A$}} $J$ un idéal de type fini dans $B=A[T]$, sa trace $I$ sur $A$ est-elle de type fini ??

% lines 18–20
[ Plus gén, est-il vrai que les $I_{\unsure{m|n}}\subset B_m$ sont des $A$-modules de type fini ?? ]

% lines bottom right
\margin{OK. ds le \textquotedblleft cas classique\textquotedblright{} $X$ \deleted{\unsure{un}} \unsure{polyèdres|polyèdre|polycylindre} \unsure{analytiques|analytique} \ldots}
